% SPDX-License-Identifier: MIT
% Copyright (c) 2026 Evgeny Baulin
% Optimization Methods, Seminar 05: theory handout (English).
% Generated by checks/05. Second-Order and Quasi-Newton Methods.ipynb from web/05/js/content-*.js and content/students.json. Do not edit by hand.
% Compile with: latexmk -xelatex -interaction=nonstopmode -halt-on-error Theory_en.tex
\documentclass[11pt,a4paper]{article}

\usepackage{fontspec}
\setmainfont{cmunrm.otf}[BoldFont=cmunbx.otf, ItalicFont=cmunti.otf, BoldItalicFont=cmunbi.otf]
\setsansfont{cmunss.otf}[BoldFont=cmunsx.otf, ItalicFont=cmunsi.otf, BoldItalicFont=cmunso.otf]
\setmonofont{cmuntt.otf}[BoldFont=cmuntb.otf, ItalicFont=cmunit.otf, BoldItalicFont=cmuntx.otf]
\usepackage{polyglossia}
\setdefaultlanguage{english}

\usepackage[a4paper,left=2.3cm,right=2.3cm,top=2.4cm,bottom=2.4cm,headheight=14pt]{geometry}
\usepackage{amsmath,amssymb}
\usepackage{graphicx}
\usepackage{array,booktabs,longtable}
\usepackage{enumitem}
\usepackage{xcolor}
\usepackage[most]{tcolorbox}
\usepackage{titlesec}
\usepackage{fancyhdr}
\usepackage{needspace}
\usepackage[expansion=false]{microtype}
\usepackage[colorlinks=true,linkcolor=hseblue,urlcolor=hseblue,citecolor=hseblue]{hyperref}

\definecolor{hseblue}{RGB}{16,45,105}
\definecolor{hseaccent}{RGB}{198,40,40}
\definecolor{hsegray}{RGB}{108,114,126}
\definecolor{hselight}{RGB}{240,243,248}
\definecolor{hsewarm}{RGB}{253,244,232}

\hypersetup{pdftitle={Optimization Methods. Seminar 05. Theory handout}, pdfauthor={Evgeny Baulin}}

\titleformat{\section}{\color{hseblue}\Large\bfseries}{\thesection}{0.6em}{}
\titleformat{\subsection}{\color{hseblue}\large\bfseries}{\thesubsection}{0.6em}{}
\titlespacing*{\section}{0pt}{16pt plus 4pt minus 2pt}{8pt plus 2pt}
\titlespacing*{\subsection}{0pt}{12pt plus 3pt minus 2pt}{6pt plus 2pt}
% a heading keeps a few lines under it on its page: a breakable box right after it would otherwise leave it alone
\pretocmd{\section}{\Needspace{7\baselineskip}}{}{\errmessage{cannot patch section}}
\pretocmd{\subsection}{\Needspace{5\baselineskip}}{}{\errmessage{cannot patch subsection}}

\tcbset{handoutbox/.style={enhanced, breakable, sharp corners, boxrule=0pt, left=7pt, right=7pt, top=5pt, bottom=5pt, before skip=8pt plus 2pt, after skip=8pt plus 2pt}}
\newtcolorbox{examplebox}{handoutbox, colback=hselight, colframe=hselight, before upper={\textcolor{hseblue}{\bfseries Example.}\ }}
\newtcolorbox{whybox}{handoutbox, colback=white, colframe=hseblue, leftrule=2.5pt, top=3pt, bottom=3pt, before upper={\textcolor{hseaccent}{\bfseries Why.}\ }}
\newtcolorbox{answerbox}{handoutbox, colback=hsewarm, colframe=hsewarm, before upper={\textcolor{hseaccent}{\bfseries Answer.}\ }}
\newtcolorbox{conventionbox}{handoutbox, colback=hselight, colframe=hseblue, leftrule=2.5pt, before upper={\textcolor{hseblue}{\bfseries Notation.}\par\smallskip}}

\newcommand{\handoutfigure}[3]{\par\medskip\noindent\begin{minipage}{\linewidth}\centering
  \includegraphics[width=#1\linewidth]{figures/#2.pdf}\par\smallskip
  \begin{minipage}{0.94\linewidth}\small\textcolor{hseblue}{\bfseries Figure.}\ #3\end{minipage}\end{minipage}\par\medskip}
\newcommand{\problemhead}[3]{\par\bigskip\noindent{\color{hseblue}\large\bfseries #1.\ #2}\par\nopagebreak
  \noindent{\small\itshape\color{hsegray}#3}\par\nopagebreak\smallskip}
\newtcolorbox{definitionbox}[1]{handoutbox, colback=white, colframe=hseblue, leftrule=2.5pt, before upper={\textcolor{hseblue}{\bfseries #1}\ }}
\newtcolorbox{theorembox}[1]{handoutbox, colback=hselight, colframe=hseblue, leftrule=2.5pt, before upper={\textcolor{hseblue}{\bfseries #1}\ }}
\newtcolorbox{proofbox}{handoutbox, colback=white, colframe=hsegray, leftrule=1.2pt, top=3pt, bottom=3pt, before upper={\textcolor{hsegray}{\bfseries\itshape Proof.}\ }}
\newtcolorbox{notebox}[1]{handoutbox, colback=hselight, colframe=hselight, before upper={\textcolor{hseblue}{\bfseries #1}\par\smallskip}}
\newtcolorbox{trapbox}{handoutbox, colback=hsewarm, colframe=hseaccent, leftrule=2.5pt, before upper={\textcolor{hseaccent}{\bfseries Trap.}\ }}
\newtcolorbox{selfcheckbox}{handoutbox, colback=white, colframe=hseaccent, leftrule=2.5pt, top=3pt, bottom=3pt, before upper={\textcolor{hseaccent}{\bfseries Check yourself.}\ }}
% a display wider than the line is set as a centred paragraph that may break after a relation
\newcommand{\handoutdisplay}[1]{\sbox0{$\displaystyle #1$}\ifdim\wd0>\linewidth
  \par{\centering\binoppenalty=10000\relpenalty=0\relax$\displaystyle #1$\par}\else\[#1\]\fi}
\newcommand{\selfcheckanswer}{\par\smallskip\textcolor{hsegray}{\bfseries Answer.}\ }

\setlist[itemize]{itemsep=3pt,topsep=4pt}
\setlist[enumerate]{itemsep=3pt,topsep=4pt}
\setlength{\parindent}{0pt}
\setlength{\parskip}{5pt plus 1pt}
\setlength{\emergencystretch}{2em}
\setcounter{tocdepth}{2}

\pagestyle{fancy}
\fancyhf{}
\fancyhead[L]{\small\color{hsegray}Optimization Methods}
\fancyhead[R]{\small\color{hsegray}Seminar 05}
\fancyfoot[L]{\small\color{hsegray}Evgeny Baulin}
\fancyfoot[C]{\small\thepage}
\fancyfoot[R]{\small\color{hsegray}MIT License}
\renewcommand{\headrulewidth}{0pt}
\renewcommand{\footrulewidth}{0pt}
\fancypagestyle{plain}{\pagestyle{fancy}}

\begin{document}

\begin{center}
  {\large\color{hsegray}Optimization Methods\par}
  \vspace{4pt}
  {\large\bfseries\color{hseaccent}Seminar 05\par}
  \vspace{8pt}
  {\LARGE\bfseries\color{hseblue}Second-Order and Quasi-Newton Methods\par}
  \vspace{8pt}
  {\Large Theory handout\par}
  \vspace{12pt}
  HSE University, Faculty of Computer Science\par
  Academic year 2026–2027\par
  \vspace{8pt}
  Course lecturer: Andrey Ignatov\par
  Evgeny Baulin\par
  Telegram: \href{https://t.me/tarakan_tuc}{@tarakan\_tuc}\par
\end{center}

\noindent{\color{hseblue}\rule{\linewidth}{0.6pt}}

\tableofcontents

\clearpage


Gradient descent uses only the slope and pays for it with the condition number. This handout follows the lecture and the computer seminar of Andrey Ignatov on second-order methods: the Newton step and its local theory, the safeguards that make it work far from the solution, Gauss–Newton and Levenberg–Marquardt, BFGS and L-BFGS, and the cost of curvature.

\begin{conventionbox}
$f:\mathbb R^n\to\mathbb R$ is twice continuously differentiable, $g=\nabla f(x)$, $H=\nabla^2 f(x)$, $e=x-x^*$. $B$ approximates the Hessian and $G$ its inverse; $s_k=x_{k+1}-x_k$, $y_k=g_{k+1}-g_k$. $A\succ0$ means positive definite. Armijo's constant is $c\in(0,1)$, usually small.
\end{conventionbox}

% =====================================================================
\section{The quadratic model and the Newton step}\label{sec:model}

\begin{definitionbox}{Definition.}
The \textbf{quadratic model} of $f$ at $x$ is
\handoutdisplay{
m_x(p)=f(x)+g^\top p+\tfrac{1}{2}p^\top Hp,
}
 which agrees with $f(x+p)$ up to $o(\|p\|^2)$. If $H$ is invertible, the \textbf{Newton direction} $p_N$ solves $Hp_N=-g$, and the full Newton step is $x^+=x+p_N$.
\end{definitionbox}

$\nabla_p m_x(p)=g+Hp$, so $p_N$ is the stationary point of the model. When $H\succ0$ it is the unique minimizer of the model; when $H$ is indefinite it is a saddle of the model, and when $H\prec0$ its maximizer.

\begin{whybox}
Gradient descent minimizes the model with $H$ replaced by $\alpha^{-1}I$: a round bowl. Newton uses the bowl of the true shape. In the eigenbasis of $H\succ0$ the Newton step divides each gradient component by its own curvature, so steep directions get short steps and flat directions long ones.
\end{whybox}

\begin{theorembox}{Theorem: One step on a quadratic.}
For $f(x)=\tfrac{1}{2}x^\top Qx-b^\top x+c$ with $Q\succ0$, one full Newton step from any $x$ reaches $x^*=Q^{-1}b$.
\end{theorembox}

\begin{proofbox}
$g=Qx-b$ and $H=Q$. The Newton equation $Qp=-(Qx-b)$ gives $p=-x+Q^{-1}b=x^*-x$, so $x+p=x^*$.
\end{proofbox}

\begin{examplebox}
For $Q=\operatorname{diag}(1,\ 100)$, GD with $\alpha=1/L=0.01$ shrinks the error along the flat axis by $0.99$ per step and needs $1375$ steps to reduce it by a factor of $1000000$. Newton's step divides the two gradient components by their curvatures and finishes in one step.
\end{examplebox}

\begin{trapbox}
Solve $Hp=-g$ with a Cholesky or LU factorization, \texttt{numpy.linalg.solve}, never as \texttt{inv(H) @ g}. The explicit inverse costs more and is less accurate, and a failed Cholesky factorization is the cheapest test that $H\succ0$.
\end{trapbox}

\problemhead{A2}{One step on a quadratic}{On your own}
\textbf{Statement.} Let $f(x)=\tfrac{1}{2}x^\top Qx-b^\top x$ with $Q=\begin{pmatrix} 4 & 1 \\ 1 & 2 \end{pmatrix}$, $b=(1,\ -1)$. From $x_0=(2,\ 0)$ compute the Newton step and $x_1$, and show that $x_1$ does not depend on $x_0$. Find the condition number $\kappa$ of $Q$.

\textbf{Solution.}
\begingroup\sloppy
\begin{enumerate}[label=\textbf{\arabic*.}, leftmargin=*]
  \item $g_0=Qx_0-b=(7,\ 3)$. Solving $Qp=-g_0$ gives $p_0=(-1.5714,\ -0.7143)$ and $x_1=x_0+p_0=(0.4286,\ -0.7143)$.
  \item For any $x$ the Newton equation reads $Qp=-(Qx-b)$, so $p=Q^{-1}b-x=x^*-x$ and $x+p=x^*$. Indeed $Qx_1-b=0$.
  \item The eigenvalues of $Q$ are $(1.5858,\ 4.4142)$, so $\kappa=2.7836$. Gradient descent with its best fixed step would contract the error by $(\kappa-1)/(\kappa+1)=0.4714$ per step.
\end{enumerate}
\endgroup
\begin{answerbox}
$x_1=(0.4286,\ -0.7143)=Q^{-1}b$ from any start; $\kappa=2.7836$.
\end{answerbox}

% =====================================================================
\section{Local quadratic convergence}\label{sec:local}

\begin{theorembox}{Theorem: Quadratic convergence.}
Let $\nabla f(x^*)=0$, $\nabla^2 f(x^*)\succ0$, and let the Hessian be $M$-Lipschitz near $x^*$. Then there are a neighbourhood of $x^*$ and $m>0$ such that full Newton steps from it stay in it and satisfy
\handoutdisplay{
\|e^+\|\le{M\over 2m}\|e\|^2.
}

\end{theorembox}

\begin{proofbox}
By the fundamental theorem of calculus, $g(x)=\int_0^1 H(x^*+te)\,e\,dt$. Since $e^+=e-H(x)^{-1}g(x)$,
\handoutdisplay{
e^+=H(x)^{-1}\int_0^1\bigl[H(x)-H(x^*+te)\bigr]e\,dt.
}
 Near $x^*$ the Hessian stays positive definite, $\|H(x)^{-1}\|\le1/m$, and $\|H(x)-H(x^*+te)\|\le M(1-t)\|e\|$. Integrating $1-t$ gives the factor $\tfrac{1}{2}$.
\end{proofbox}

Once $\tfrac{M}{2m}\|e\|<1$, each step about doubles how many leading digits of $x_k$ are right. Both constants matter: a nearly singular Hessian at the solution (small $m$) or a fast-changing one (large $M$) shrinks the region where this happens.

\handoutfigure{0.85}{fig_digits}{Error of Newton's method on $x-\ln x$ against the iteration, compared with gradient descent.}

\problemhead{A3}{Digits double}{On your own}
\textbf{Statement.} Apply Newton's method to $f(x)=x-\ln x$, $x>0$, from $x_0=0.5$. Derive the recurrence for $e_k=1-x_k$ and compute $x_3$. How many steps make $e_k\le10^{-12}$? For which $x_0>0$ does the method converge?

\textbf{Solution.}
\begingroup\sloppy
\begin{enumerate}[label=\textbf{\arabic*.}, leftmargin=*]
  \item $x^+=x-(1-1/x)x^2=2x-x^2$. With $e=1-x$: $e^+=1-2x+x^2=(1-x)^2=e^2$, exactly, with $C=1$.
  \item From $x_0=0.5$ the iterates are $(0.5,\ 0.75,\ 0.9375,\ 0.9961)$, the errors $(0.5,\ 0.25,\ 0.0625,\ 0.0039)$; $x_3=0.9961$.
  \item $e_k=e_0^{2^k}=2^{-2^k}$. It falls below $10^{-12}$ first at $k=6$: every step doubles the number of correct binary digits.
  \item The iteration converges exactly when $|e_0|<1$, that is $0< x_0<2$. From $x_0=2$ it lands on $0$, outside the domain; from a larger start, for example $x_0=2.5$, it jumps to $-1.25$.
\end{enumerate}
\endgroup
\begin{answerbox}
$e^+=e^2$, $x_3=0.9961$, $6$ steps; convergence for $0< x_0<2$.
\end{answerbox}

\begin{trapbox}
For $f(x)=x^4$ the minimizer has $f''(0)=0$, and Newton gives $x^+=0.6667\,x$: linear convergence, $35$ steps per factor of $1000000$. The theorem above assumes $\nabla^2 f(x^*)\succ0$ for a reason.
\end{trapbox}

% =====================================================================
\section{When the full step fails}\label{sec:failure}

Along $p_N$, $g^\top p_N=-g^\top H^{-1}g$. This is negative when $H\succ0$ and can be positive when $H$ is indefinite. Even with $H\succ0$ the unit step can leave the region where the model describes $f$.

\begin{itemize}
  \item \textbf{Uphill direction.} An indefinite $H$ may give $g^\top p_N>0$; a line search along $p_N$ cannot help.
  \item \textbf{Wrong stationary point.} Newton solves $\nabla f=0$ and is attracted by saddles and maxima as readily as by minima.
  \item \textbf{Overlong step.} The model is right at $x$ and wrong where its minimizer lies.
  \item **Singular or ill-conditioned $H$.** The linear solve amplifies rounding errors, and the step may be huge.
\end{itemize}

\problemhead{B2}{Attracted to a maximum}{On your own}
\textbf{Statement.} Apply full Newton to $f(x)=\tfrac{1}{4}x^4-\tfrac{1}{2}x^2$ from $x_0=0.4$. Write the iteration as $x^+=\varphi(x)$ and compute $x_1$, $x_2$. Where does it go, and what kind of point is that? Find the largest $a$ such that every start with $0<|x_0|< a$ is drawn to that point.

\textbf{Solution.}
\begingroup\sloppy
\begin{enumerate}[label=\textbf{\arabic*.}, leftmargin=*]
  \item $\varphi(x)=x-{x^3-x\over 3x^2-1}={2x^3\over 3x^2-1}$. From $x_0=0.4$: $x_1=-0.2462$, $x_2=0.0365$.
  \item The iterates go to $0$, where $f'(0)=0$ and $f''(0)=-1$: a local maximum. Newton looks for $f'=0$ and does not tell a maximum from a minimum; near $0$, $\varphi(x)\approx-2x^3$, so the convergence is even cubic.
  \item $\varphi(x)=-x$ gives $2x^2=1-3x^2$, $x^2=1/5$. For $0<|x|< a$ with $a=1/\sqrt5=0.4472$ one checks $|\varphi(x)|<|x|$, so the iterates shrink to the maximum; at $|x|=a$ they cycle.
\end{enumerate}
\endgroup
\begin{answerbox}
$x_1=-0.2462$, $x_2=0.0365$; convergence to the maximum $x=0$ for $0<|x_0|<0.4472$.
\end{answerbox}

\begin{definitionbox}{Definition.}
\textbf{Damped Newton} chooses $\alpha_k=1,\beta,\beta^2,\dots$ until
\handoutdisplay{
f(x_k+\alpha_kp_k)\le f(x_k)+c\,\alpha_k\nabla f(x_k)^\top p_k,
}
 and sets $x_{k+1}=x_k+\alpha_kp_k$. It needs $p_k$ to be a descent direction.
\end{definitionbox}

\begin{theorembox}{Proposition.}
If the iterates of damped Newton converge to $x^*$ with $\nabla^2 f(x^*)\succ0$ and $c<\tfrac{1}{2}$, then $\alpha_k=1$ is accepted for all large $k$, and the convergence is quadratic.
\end{theorembox}

\begin{proofbox}
Near $x^*$, $f(x+p_N)-f(x)=-\tfrac{1}{2}g^\top H^{-1}g+o(\|g\|^2)$, while the Armijo bound asks for at most $-c\,g^\top H^{-1}g$. Since $c<\tfrac{1}{2}$, the first is eventually below the second.
\end{proofbox}

\problemhead{A1}{Newton on Rosenbrock's function: a step too far}{On the board}
\textbf{Statement.} For $f(x,y)=a(y-x^2)^2+(1-x)^2$ with $a=100$, start at $x_0=(0,\ 0)$. Take two full Newton steps and compute $f$ after each. Then run damped Newton from $x_0$: halve $\alpha$ from $1$ until the Armijo condition with $c=0.0001$ holds. Which $\alpha$ is accepted, and what is $f$ there?

\textbf{Solution.}
\begingroup\sloppy
\begin{enumerate}[label=\textbf{\arabic*.}, leftmargin=*]
  \item At $x_0$: $\nabla f=(-2,\ 0)$ and $\nabla^2 f=\begin{pmatrix} 2 & 0 \\ 0 & 200 \end{pmatrix}$. Solving $\nabla^2 f\,p=-\nabla f$ gives $p_0=(1,\ 0)$, so $x_1=(1,\ 0)$.
  \item But $f(x_1)=100$, while $f(x_0)=1$. The quadratic model at $x_0$ is exact up to second order, and its minimizer lies far outside the region where it describes $f$: at $x=1$ the valley $y=x^2$ is at height one, and the step stays at $y=0$.
  \item From $x_1$: $\nabla f=(400,\ -200)$, $\nabla^2 f=\begin{pmatrix} 1202 & -400 \\ -400 & 200 \end{pmatrix}$, and the Newton step is $p_1=(0,\ 1)$. So $x_2=(1,\ 1)$ is the minimizer, $f(x_2)=0$. The full method gets there, but through a point a hundred times worse than the start.
  \item Damped Newton keeps the direction $p_0$ and tests $f(x_0+\alpha p_0)\le f(x_0)+c\alpha\nabla f^\top p_0$, where $\nabla f^\top p_0=-2$. For $\alpha=(1,\ 0.5,\ 0.25)$ the values are $f=(100,\ 6.5,\ 0.9531)$: the first two fail, and $\alpha=0.25$ is accepted at $x=(0.25,\ 0)$.
\end{enumerate}
\endgroup
\begin{answerbox}
$f(x_1)=100$, $x_2=(1,\ 1)$; damped: $\alpha=0.25$, $f=(100,\ 6.5,\ 0.9531)$ for the trials, the last accepted.
\end{answerbox}

% =====================================================================
\section{Shifting the Hessian}\label{sec:regularization}

\begin{definitionbox}{Definition.}
The \textbf{regularized Newton step} solves $(H+\lambda I)p=-g$ with $\lambda\ge0$ chosen so that $H+\lambda I\succ0$, that is $\lambda>-\lambda_{\min}(H)$.
\end{definitionbox}

\begin{theorembox}{Proposition.}
If $H+\lambda I\succ0$ and $g\ne0$, then $g^\top p=-g^\top(H+\lambda I)^{-1}g<0$. As $\lambda\to0$ the step tends to the Newton step when $H\succ0$; as $\lambda\to\infty$, $p\approx-\lambda^{-1}g$, a short gradient step.
\end{theorembox}

In practice $\lambda$ starts at zero or at a small value and is multiplied by a constant until the Cholesky factorization of $H+\lambda I$ succeeds; after a successful full step it is reduced again. In the eigenbasis the step is $p_i=-g_i/(\lambda_i+\lambda)$: a negative-curvature direction with $\lambda_i+\lambda>0$ is now followed downhill.

\problemhead{B1}{An uphill Newton direction}{On the board}
\textbf{Statement.} Let $f(x,y)=\tfrac{1}{2}x^2-y^2+\tfrac{1}{4}y^4$ at the point $(0,\ 0.5)$. Compute $g$, $H$, the Newton direction $p_N$ and $g^\top p_N$. Is $p_N$ a descent direction? Find the smallest shift that makes $H+\lambda I$ positive definite, solve $(H+\lambda I)p=-g$ for $\lambda=2$ and test the full step with Armijo, $c=0.0001$.

\textbf{Solution.}
\begingroup\sloppy
\begin{enumerate}[label=\textbf{\arabic*.}, leftmargin=*]
  \item $g=(0,\ -0.875)$, $H=\begin{pmatrix} 1 & 0 \\ 0 & -1.25 \end{pmatrix}$: one eigenvalue is negative, so $H$ is indefinite. Newton gives $p_N=-H^{-1}g=(0,\ -0.7)$ and $g^\top p_N=0.6125>0$.
  \item $\phi(\alpha)=f(x+\alpha p_N)$ has $\phi'(0)=g^\top p_N>0$: every small step raises $f$, and the full step takes it from $-0.2344$ to $-0.0396$. Backtracking would shrink $\alpha$ forever. The direction itself must change.
  \item $H+\lambda I\succ0$ exactly when $\lambda>1.25$. For $\lambda=2$: $p=(0,\ 1.1667)$ and $g^\top p=-1.0208<0$. The step now goes along the negative curvature, away from the saddle.
  \item Armijo at $\alpha=1$: $f(x+p)=-0.8488\le f(x)+c\,g^\top p=-0.2345$, so the full regularized step is accepted. The minimizers are at $y=\pm1.4142$, where $H\succ0$ and the shift can be dropped.
\end{enumerate}
\endgroup
\begin{answerbox}
$g^\top p_N=0.6125$: not a descent direction. $\lambda>1.25$; for $\lambda=2$ the step is $(0,\ 1.1667)$ and $f$ falls to $-0.8488$.
\end{answerbox}

\handoutfigure{0.95}{fig_saddle}{Full Newton converges to the saddle; the regularized, damped method leaves it along the negative curvature.}

\begin{notebox}{Trust regions}
Trust-region methods choose the step size first and the shift second: $p$ minimizes the model over $\|p\|\le\Delta$, and the solution has the form $(H+\lambda I)p=-g$ with $\lambda\ge0$ fixed by $\Delta$. They handle negative curvature without a separate test.
\end{notebox}

% =====================================================================
\section{Gauss–Newton and Levenberg–Marquardt}\label{sec:least-squares}

For $f(x)=\tfrac{1}{2}\|r(x)\|^2$ with residuals $r:\mathbb R^n\to\mathbb R^N$ and Jacobian $J$,
\handoutdisplay{
\nabla f=J^\top r,\qquad \nabla^2 f=J^\top J+\sum_i r_i\nabla^2 r_i.
}


\begin{definitionbox}{Definition.}
\textbf{Gauss–Newton} drops the second term and solves $J^\top Jp=-J^\top r$; it is the Newton step for the linearized residual $r+Jp$. \textbf{Levenberg–Marquardt} solves $(J^\top J+\lambda I)p=-J^\top r$.
\end{definitionbox}

\begin{whybox}
$J^\top J\succeq0$ needs only first derivatives, and the dropped term is small when the residuals are small near the solution. The shift $\lambda$ handles a rank-deficient $J$ and moves the step from Gauss–Newton toward a short gradient step.
\end{whybox}

\problemhead{B3}{Gauss–Newton, Newton and Levenberg–Marquardt}{On your own}
\textbf{Statement.} Let $f(\theta)=\tfrac{1}{2}\|r(\theta)\|^2$ with $r(\theta)=(\theta^2-2,\ \theta-1)$. From $\theta_0=1$ take one Gauss–Newton step, one Newton step with the exact second derivative and one Levenberg–Marquardt step with $\lambda=1$. Report $\theta_1$ for each and compare with the minimizer.

\textbf{Solution.}
\begingroup\sloppy
\begin{enumerate}[label=\textbf{\arabic*.}, leftmargin=*]
  \item At $\theta_0$: $r=(-1,\ 0)$, $J=(2,\ 1)$, $J^\top J=5$, $J^\top r=-2$. Gauss–Newton: $\theta_1=\theta_0-J^\top r/J^\top J=1.4$.
  \item The exact second derivative adds $r_1r_1''=-2$, so $f''=3$ and Newton gives $\theta_1=1.6667$. Levenberg–Marquardt divides by $J^\top J+\lambda$: $\theta_1=1.3333$.
  \item $f'(\theta)=2\theta^3-3\theta-1=(\theta+1)(2\theta^2-2\theta-1)$, and the minimizer is $\theta^*=(1+\sqrt3)/2=1.366$. Gauss–Newton lands closest: at $\theta_0$ the residual term made the exact curvature smaller and the Newton step too long.
\end{enumerate}
\endgroup
\begin{answerbox}
Gauss–Newton $1.4$, Newton $1.6667$, Levenberg–Marquardt $1.3333$; $\theta^*=1.366$.
\end{answerbox}

% =====================================================================
\section{The cost of curvature}\label{sec:cost}

\begin{center}\small
\begin{tabular}{@{}>{\raggedright\arraybackslash\hspace{0pt}}p{0.26\linewidth}>{\raggedright\arraybackslash\hspace{0pt}}p{0.49\linewidth}l@{}}
\toprule
Operation & Dense cost & Memory \\
\midrule
Form the Hessian & $O(n^2)$ derivatives; $O(Nn^2)$ for a GLM & $O(n^2)$ \\
Cholesky factorization & $\approx n^3/3$ & $O(n^2)$ \\
Solve with the factor & $O(n^2)$ & $O(n^2)$ \\
Hessian-vector product & a few gradients & $O(n)$ extra \\
L-BFGS direction & $\approx 4mn$ & $2mn$ \\
\bottomrule
\end{tabular}
\end{center}

\begin{examplebox}
With an $\ell_2$ penalty, the logistic loss $f(w)=N^{-1}\sum_i\bigl[\log(1+e^{z_i})-y_iz_i\bigr]+\tfrac{\gamma}{2}\|w\|^2$, $z=Xw$:
\handoutdisplay{
\nabla f=N^{-1}X^\top(\sigma(z)-y)+\gamma w,\quad Hv=N^{-1}X^\top\bigl(d\odot Xv\bigr)+\gamma v,
}
 with $d_i=\sigma(z_i)(1-\sigma(z_i))$. The product costs two passes over $X$ and never forms $H$.
\end{examplebox}

\textbf{Newton-CG} solves $Hp=-g$ approximately by conjugate gradients, which need only products $Hv$. Stopping CG early, or at negative curvature $v^\top Hv\le0$, still gives a descent direction.

\problemhead{C3}{What a Hessian costs}{On your own}
\textbf{Statement.} Logistic regression with $N=100000$ samples and $n=20000$ features. Forming $H=N^{-1}X^\top DX$ costs about $Nn^2$ multiply-adds; one product $Hv=N^{-1}X^\top(d\odot Xv)$ costs about $2Nn$. How many times cheaper is a Newton-CG solve with $k=16$ products than forming $H$? How many gigabytes does a dense $H$ take with $8$ bytes per entry, and how many megabytes do $m=5$ L-BFGS pairs take?

\textbf{Solution.}
\begingroup\sloppy
\begin{enumerate}[label=\textbf{\arabic*.}, leftmargin=*]
  \item $Nn^2/(2kNn)=n/(2k)=625$: the samples cancel, and the gain grows with the number of features.
  \item A dense $H$ has $n^2$ entries: $3.2$ GB, before any factorization. L-BFGS stores $2m$ vectors of length $n$: $1.6$ MB.
  \item A Cholesky factorization would add about $n^3/3$ operations. With $n$ in the tens of thousands, Newton is possible only matrix-free, through $Hv$ and conjugate gradients.
\end{enumerate}
\endgroup
\begin{answerbox}
Ratio $625$; dense $H$ $3.2$ GB; L-BFGS $1.6$ MB.
\end{answerbox}

\handoutfigure{0.95}{fig_cost}{Gradient norm against iterations and against passes over the data for gradient descent, damped Newton, BFGS and L-BFGS.}

% =====================================================================
\section{Secant pairs and BFGS}\label{sec:bfgs}

By the fundamental theorem of calculus,
\handoutdisplay{
y_k=\Bigl[\int_0^1\nabla^2 f(x_k+ts_k)\,dt\Bigr]s_k:
}
 $y_k$ is the action of an average Hessian on $s_k$. A quasi-Newton method keeps a matrix that reproduces it.

\begin{definitionbox}{Definition.}
The \textbf{secant conditions} are $B_{k+1}s_k=y_k$ and $G_{k+1}y_k=s_k$. The \textbf{BFGS update}, in Hessian and inverse form, is
\handoutdisplay{
B_{k+1}=B_k-{B_ks_ks_k^\top B_k\over s_k^\top B_ks_k}+{y_ky_k^\top\over y_k^\top s_k},\qquad G_{k+1}=(I-\rho_ks_ky_k^\top)G_k(I-\rho_ky_ks_k^\top)+\rho_ks_ks_k^\top,
}
 with $\rho_k=1/(y_k^\top s_k)$. The direction is $p_k=-G_kg_k$, with a line search.
\end{definitionbox}

\begin{notebox}{Why this formula}
One vector equation does not determine a matrix. BFGS is the symmetric update closest to $G_k$ in a weighted norm that satisfies the secant condition; it changes $G$ by a matrix of rank at most two. DFP is the same construction with the roles of $B$ and $G$ exchanged.
\end{notebox}

\begin{theorembox}{Theorem.}
$G_{k+1}y_k=s_k$. If $G_k\succ0$ and $y_k^\top s_k>0$, then $G_{k+1}\succ0$.
\end{theorembox}

\begin{proofbox}
$(I-\rho ys^\top)y=y-\rho y(s^\top y)=0$, so the first term vanishes on $y$, and $\rho ss^\top y=s$. For $z\ne0$ put $V=I-\rho sy^\top$: $z^\top G_{k+1}z=(V^\top z)^\top G_k(V^\top z)+\rho(s^\top z)^2\ge0$. If it were zero, then $s^\top z=0$, hence $V^\top z=z-\rho y(s^\top z)=z$, and $z^\top G_kz=0$, so $z=0$.
\end{proofbox}

\problemhead{C1}{BFGS on a quadratic, by hand}{On the board}
\textbf{Statement.} Let $f(x)=\tfrac{1}{2}x^\top Qx$ with $Q=\begin{pmatrix} 2 & 1 \\ 1 & 1 \end{pmatrix}$, $x_0=(-2,\ 2)$, $G_0=I$, and use exact line searches. Take the first step, form $s$, $y$ and the inverse BFGS update $G_1$. Check the secant condition and positive definiteness. Take the second step: what are $\alpha_1$ and $x_2$?

\textbf{Solution.}
\begingroup\sloppy
\begin{enumerate}[label=\textbf{\arabic*.}, leftmargin=*]
  \item $g_0=Qx_0=(-2,\ 0)$, $p_0=-G_0g_0=(2,\ 0)$, $\alpha_0=0.5$. So $s=(1,\ 0)$, $x_1=(-1,\ 2)$, $g_1=(0,\ 1)$ and $y=g_1-g_0=(2,\ 1)$; on a quadratic $y=Qs$ exactly.
  \item $y^\top s=2>0$, $\rho=0.5$. With $V=I-\rho sy^\top$,
\handoutdisplay{
G_1=VG_0V^\top+\rho ss^\top=\begin{pmatrix} 0.75 & -0.5 \\ -0.5 & 1 \end{pmatrix}.
}

  \item $G_1y=(1,\ 0)=s$. $G_1$ is symmetric with eigenvalues $(0.3596,\ 1.3904)$, both positive. Its inverse $B_1=\begin{pmatrix} 2 & 1 \\ 1 & 1.5 \end{pmatrix}$ agrees with $Q$ along $s$ and nowhere else yet.
  \item $p_1=-G_1g_1=(0.5,\ -1)$, $g_1^\top p_1<0$, and the exact step is $\alpha_1=2$: $x_2=(0,\ 0)=x^*$. With exact line searches BFGS minimizes an $n$-dimensional positive definite quadratic in at most $n$ steps.
\end{enumerate}
\endgroup
\begin{answerbox}
$G_1=\begin{pmatrix} 0.75 & -0.5 \\ -0.5 & 1 \end{pmatrix}$, $G_1y=s$; $\alpha_1=2$, $x_2=(0,\ 0)$.
\end{answerbox}

\handoutfigure{0.85}{fig_bfgs}{Gradient descent, BFGS and Newton on Rosenbrock's function: objective gap against the iteration.}

Near a regular minimizer BFGS with a line search that eventually accepts $\alpha=1$ converges superlinearly: $\|e_{k+1}\|/\|e_k\|\to0$. It does not reach Newton's quadratic rate, but it needs only gradients and $O(n^2)$ work per step.

% =====================================================================
\section{The curvature condition}\label{sec:curvature}

$y_k^\top s_k=s_k^\top\bigl[\int_0^1\nabla^2 f(x_k+ts_k)dt\bigr]s_k$ is the average curvature along the step. For a $\mu$-strongly convex $f$ it is at least $\mu\|s_k\|^2>0$; for a nonconvex $f$ it can have either sign.

\begin{definitionbox}{Definition.}
The \textbf{Wolfe conditions} for a step $\alpha$ along a descent direction $p$ are the Armijo condition and the curvature condition $\nabla f(x+\alpha p)^\top p\ge c_2\nabla f(x)^\top p$, with $0< c< c_2<1$.
\end{definitionbox}

\begin{theorembox}{Proposition.}
If $\alpha$ satisfies the Wolfe curvature condition along a descent direction $p$, then $y^\top s>0$ for $s=\alpha p$.
\end{theorembox}

\begin{proofbox}
$y^\top s=\alpha\bigl(\nabla f(x+\alpha p)-\nabla f(x)\bigr)^\top p\ge\alpha(c_2-1)\nabla f(x)^\top p>0$, because $c_2<1$ and $\nabla f(x)^\top p<0$.
\end{proofbox}

\problemhead{C2}{The curvature condition}{On your own}
\textbf{Statement.} For $f(x)=\tfrac{1}{4}x^4-\tfrac{1}{2}x^2$ take the points $(0,\ 0.5,\ 1.5)$. Compute $y^\top s$ for the step from the first point to the second; may BFGS use this pair? For the step from the second to the third compute the secant curvature $B=y/s$ and the quasi-Newton step $x^+=x-f'(x)/B$ from the third point.

\textbf{Solution.}
\begingroup\sloppy
\begin{enumerate}[label=\textbf{\arabic*.}, leftmargin=*]
  \item From the first point to the second: $s=0.5$, $y=-0.375$ and $y^\top s=-0.1875<0$. The update would have to make $G$ negative; BFGS must skip this pair.
  \item From the second to the third: $s=1$, $y=2.25$, $B=y/s=2.25$. By the mean value theorem this is the average of $f''$ over the step, and it is positive.
  \item From the third point, $x^+=x-f'(x)/B=0.6667$. The minimizer is $1$; in one dimension BFGS is the secant method, which converges superlinearly near a regular minimizer.
\end{enumerate}
\endgroup
\begin{answerbox}
$y^\top s=-0.1875$: skip the pair. $B=2.25$, $x^+=0.6667$.
\end{answerbox}

\begin{trapbox}
Armijo backtracking alone does not guarantee $y^\top s>0$. Implementations skip the update when $y^\top s\le\varepsilon\|y\|\|s\|$ or damp $y$ toward $B s$ (Powell's damping). Without the test one update with $y^\top s<0$ can make $G$ indefinite and the next direction an ascent direction.
\end{trapbox}

% =====================================================================
\section{Limited memory and bounds}\label{sec:lbfgs}

Dense BFGS stores an $n\times n$ matrix. \textbf{L-BFGS} stores the last $m$ pairs $(s_i,y_i)$, usually $m$ between five and twenty, and applies $G_k$ to $g_k$ without forming it.

\begin{itemize}
  \item Set $q=g_k$. For the pairs from the newest to the oldest: $a_i=\rho_is_i^\top q$, $q\leftarrow q-a_iy_i$.
  \item Scale: $r=\gamma_kq$ with $\gamma_k=s_{k-1}^\top y_{k-1}/y_{k-1}^\top y_{k-1}$, an estimate of the inverse curvature.
  \item For the pairs from the oldest to the newest: $b_i=\rho_iy_i^\top r$, $r\leftarrow r+s_i(a_i-b_i)$.
  \item Return $p_k=-r$: about $4mn$ operations and $2mn$ stored numbers.
\end{itemize}

The result equals $-Gg$ for the matrix obtained from $\gamma_kI$ by $m$ BFGS updates with the stored pairs. Old curvature is forgotten, which costs iterations on problems where it would still be useful.

\begin{notebox}{L-BFGS-B}
For bounds $l\le x\le u$, \textbf{L-BFGS-B} moves along the projected gradient path to a generalized Cauchy point, fixes the variables at their bounds, minimizes the limited-memory model over the free ones and runs a line search that respects the box. Stationarity is measured by the projected gradient.
\end{notebox}

\begin{trapbox}
Clipping the unconstrained L-BFGS solution to the box is not the constrained minimizer in general: when two variables are coupled, fixing one at its bound moves the optimum of the other.
\end{trapbox}

\begin{center}\small
\begin{tabular}{@{}lll>{\raggedright\arraybackslash\hspace{0pt}}p{0.45\linewidth}@{}}
\toprule
Method & Needs & Per step & Near the solution \\
\midrule
Gradient descent & $g$ & $O(n)$ & linear, rate set by $\kappa$ \\
Newton & $g$, $H$ & $O(n^3)$ & quadratic \\
Newton-CG & $g$, $Hv$ & CG iterations & superlinear with a tightening tolerance \\
BFGS & $g$ & $O(n^2)$ & superlinear \\
L-BFGS & $g$ & $O(mn)$ & linear in general \\
\bottomrule
\end{tabular}
\end{center}

% =====================================================================
\section{Implementation checklist and exercises}\label{sec:practice}

\begin{itemize}
  \item Solve, do not invert; detect $H\not\succ0$ by a failed Cholesky factorization and shift.
  \item Test $g^\top p<0$ before the line search and $y^\top s>0$ before the BFGS update.
  \item Check derivatives: central differences for $g$, a difference of gradients for $Hv$, over several step sizes.
  \item Log function, gradient and Hessian calls next to iterations, and compare methods at one stopping rule.
  \item At the end check the Hessian or the smallest curvature found: a small gradient may mean a saddle.
\end{itemize}

\begin{selfcheckbox}
Why can backtracking not repair a Newton direction with $g^\top p>0$?
\selfcheckanswer Because $f(x+\alpha p)=f(x)+\alpha g^\top p+o(\alpha)$ rises for every small $\alpha$; the Armijo test fails for all of them.
\end{selfcheckbox}

\begin{selfcheckbox}
What information does one secant pair $(s,y)$ carry?
\selfcheckanswer The action of the average Hessian along the segment on the single vector $s$, nothing about other directions.
\end{selfcheckbox}

\begin{selfcheckbox}
When is a Hessian-vector product more useful than the Hessian?
\selfcheckanswer When $n$ is large: $Hv$ costs a few gradients and $O(n)$ memory, and conjugate gradients need only such products.
\end{selfcheckbox}

\end{document}
